\subsection{From Hand-Crafted to Data-Driven Features}\label{sec:from-hand-crafted-to-data-driven-features}

Before deep learning became popular, computers did \textbf{not} learn these visual features automatically. Instead, engineers manually designed algorithms that extracted useful information from an image before classification. The complete pipeline for image classification was as follows:
\begin{equation*}
    \text{Image} \to \text{Feature Extraction} \to \text{Classifier}
\end{equation*}
Where:
\begin{itemize}
    \item The \textbf{feature extraction} was designed by humans!
    \item Only the \textbf{classifier} was learned from data.
\end{itemize}
This chapter begins by explaining why this approach eventually became insufficient, motivating CNNs.

\longline

\subsubsection{The Feature-Extraction Perspective}\label{sec:the-feature-extraction-perspective}

This first section asks a simple but fundamental question: \emph{\textbf{why can't we classify images directly from pixels?}} The answer, already hinted at in the previous section, is that \textbf{raw data usually contains much more information than is actually useful for the task}. Our goal is therefore not simply to classify images. Our goal is first to find a \textbf{better representation} of an image.

\highspace
\begin{flushleft}
    \textcolor{Green3}{\faIcon{question-circle} \textbf{What is a feature?}}
\end{flushleft}
Let's start from the beginning. A \definition{Feature} is any \textbf{quantity extracted from the original data} that \hl{helps distinguish one class from another.}

\highspace
For example, imagine distinguishing between cats and dogs. The raw image contains thousands of pixel values. However, what really matters are properties such as ears, eyes, fur texture, snout shape, body portions. These are much more informative than the extract RGB value of pixel $\left(125, 87\right)$.

\highspace
A feature is therefore \textbf{a representation that emphasizes useful information while ignoring irrelevant details}.

\highspace
\begin{flushleft}
    \textcolor{Green3}{\faIcon{question-circle} \textbf{Why raw pixels are a poor representation}}
\end{flushleft}
Suppose we have a grayscale image of size $100 \times 100$. The classifier receives a vector of $100 \times 100 = 10,000$ pixel values. Most of these numbers are \textbf{not directly meaningful}. For example, consider two photos of the same cat. The pixel values may change because of lighting, shadows, camera position, slight translations, background, image noise, or even the cat's pose.

\highspace
Although the \textbf{object is the same}, thousands of \textbf{pixel values are different}. The classifier therefore has a difficult task: \hl{it must discover from scratch that all these images represent the same object.}

\highspace
\begin{flushleft}
    \textcolor{Green3}{\faIcon{filter} \textbf{Feature Extraction}}
\end{flushleft}
Instead of using all raw pixels, we can transform the image into a more meaningful representation. Conceptually:
\begin{equation*}
    \text{Image} \to \text{Feature Extraction} \to \text{Feature Vector}
\end{equation*}
The feature extractor \textbf{removes unnecessary information} and \textbf{keeps only what is useful} for recognition.

\highspace
For example, instead of $10,000$ pixel values, we might obtain only $100$ meaningful descriptors. Those descriptors may describe edges, corners, textures, shapes, orientations, or other properties of the image, depending on the extraction method. The \hl{classifier now receives information that is much easier to separate into classes.}

\highspace
\begin{flushleft}
    \textcolor{Green3}{\faIcon{compress-arrows-alt} \textbf{Dimensionality reduction}}
\end{flushleft}
Another advantage is that feature extraction usually \textbf{reduces the number of variables}. Instead of working with $10,000$ pixel values, we may work with only a few hundred informative features.

\highspace
This has several \textbf{advantages}:
\begin{itemize}
    \item[\textcolor{Green3}{\faIcon{check}}] Less memory is required to store the data.
    \item[\textcolor{Green3}{\faIcon{check}}] Faster computation, since the classifier has fewer variables to process.
    \item[\textcolor{Green3}{\faIcon{check}}] Fewer parameters to learn, which reduces the risk of overfitting.
    \item[\textcolor{Green3}{\faIcon{check}}] Lower risk of overfitting, since the classifier has fewer variables to process.
    \item[\textcolor{Green3}{\faIcon{check}}] More robust classification, since the classifier is less sensitive to irrelevant details in the input.
\end{itemize}
Notice that \textbf{dimensionality reduction is not the primary goal}. The real goal is to produce a \textbf{better representation}. The reduction in dimensionality is often a \textbf{side effect} of removing irrelevant information.

\highspace
\begin{flushleft}
    \textcolor{Red2}{\faIcon{exclamation-triangle} \textbf{The classifier is only as good as its features}}
\end{flushleft}
This is one of the most fundamental principles in machine learning. The performance of a model is ultimately limited by the quality of its input features. If the input data does not contain informative or discriminative features, even a very powerful model will struggle to learn a good decision boundary. This idea is commonly summarized by the expression: ``\textbf{\emph{Garbage in, garbage out.}}''.

\highspace
The same principle applies to image classification. A classifier does \textbf{not invent information}. It can only use the information it receives. Imagine two situations:
\begin{itemize}
    \item \textbf{Bad features}. Suppose we want to distinguish apples from oranges, but our feature extractor only measures \emph{average brightness}. Both fruits may have similar brightness. The classifier has almost no useful information. Even a very powerful neural network will struggle to separate the two classes.
    
    \item \textbf{Good features}. Now suppose the features include color, texture, roundness, diameter. These characteristics clearly separate the two fruits. Even a relatively simple classifier can easily distinguish between apples and oranges.
\end{itemize}
As we said before, \hl{good features are often more important than a sophisticated classifier.} In fact, before CNNs became popular, much of computer vision research focused on inventing better feature extractors rather than better classifiers.

\highspace
\begin{takeawaysbox}
    \begin{itemize}
        \item A raw image is usually \textbf{not the best representation} for classification.
        \item A \textbf{feature} is a more informative description of the image.
        \item Feature extraction removes irrelevant information while preserving useful patterns.
        \item It often reduces dimensionality, making learning easier.
        \item \hl{The quality of the classifier depends heavily on the quality of the features.}
        \item Before CNNs, these features were designed manually. CNNs will instead learn them automatically from data.
    \end{itemize}
\end{takeawaysbox}